\Askhsh{36660}{4}
\begin{ylh}{2}
    \item 5.2. Αριθμητική πρόοδος
\end{ylh}
Ένας μελισσοκόμος έχει τοποθετήσει 20 κυψέλες σε μια ευθεία η οποία διέρχεται από την αποθήκη του Α. Η πρώτη κυψέλη απέχει 1 μέτρο από την αποθήκη Α, η δεύτερη 4 μέτρα από το Α, η τρίτη 7 μέτρα από το Α και γενικά κάθε επόμενη κυψέλη απέχει από την αποθήκη Α, 3 επιπλέον μέτρα, σε σχέση με την προηγούμενη κυψέλη.

\begin{alist}
\item Να αποδείξετε ότι οι αποστάσεις των κυψελών από την αποθήκη Α αποτελούν διαδοχικούς όρους αριθμητικής προόδου και να βρείτε τον ν-οστό όρο της προόδου. Τι εκφράζει ο πρώτος όρος της αριθμητικής προόδου και τι η διαφορά της;

\monades{6}
\item Σε πόση απόσταση από την αποθήκη Α είναι η 20\textsuperscript{η} κυψέλη;
\monades{6}
\item Ο μελισσοκόμος ξεκινώντας από την αποθήκη συλλέγει το μέλι, από μια κυψέλη κάθε φορά, και το μεταφέρει στην αποθήκη Α.
\begin{rlist}
\item
  Ποια είναι η απόσταση που θα διανύσει ο μελισσοκόμος για να συλλέξει το μέλι από την 3\textsuperscript{η} κυψέλη;
\monades{6}
\item
  Ποια είναι η συνολική απόσταση που θα διανύσει ο μελισσοκόμος για να συλλέξει το μέλι και από τις 20 κυψέλες;
\monades{7}
\end{rlist}
\end{alist}

